\sect{Ćwiczenia 11}

\begin{frame}{Zadanie 1}
	\begin{figure}
		\includegraphics[scale=0.08, alt={}]{./imags/E11/01.png}
		\caption{Zadanie 1. Oblicz napięcie na kondensatorze, napięcie na rezystorze i prąd w obwodzie.}
	\end{figure}	
\end{frame}

\begin{frame}{Zadanie 2}
	\begin{figure}
		\includegraphics[scale=0.08, alt={}]{./imags/E11/02.png}
		\caption{Zadanie 2. Oblicz napięcie na cewce, napięcie na rezystorze i prąd w obwodzie.}
	\end{figure}	
\end{frame}

\begin{frame}{Zadanie 3}
	\begin{figure}
		\includegraphics[scale=0.08, alt={}]{./imags/E11/03.png}
		\caption{Zadanie 3. Oblicz prąd cewki, napięcie na kondensatorze i moc wydzielaną na rezystorze.}
	\end{figure}	
\end{frame}

\begin{frame}{Zadanie 4}
	\begin{figure}
		\includegraphics[scale=0.08, alt={}]{./imags/E11/04.png}
		\caption{Zadanie 4. Oblicz prąd cewki i napięcie na kondensatorze.}
	\end{figure}	
\end{frame}

\begin{frame}{Zadanie 5}
	\begin{figure}
		\includegraphics[scale=0.08, alt={}]{./imags/E11/05.png}
		\caption{Zadanie 5. Oblicz napięcia i natężenia na wszystkich elementach obwodu.}
	\end{figure}	
\end{frame}

\begin{frame}{Zadanie 6}
	\begin{figure}
		\includegraphics[scale=0.08, alt={}]{./imags/E11/06.png}
		\caption{Zadanie 6. Oblicz napięcia i natężenia na wszystkich elementach biernych obwodu.}
	\end{figure}	
\end{frame}
